\begin{table}[ht]
\centering
\caption{Performance Evaluation Metrics}
\label{tab:evaluation_metrics}
\begin{tabular}{|p{4cm}|p{6cm}|}
\hline
Metric & Purpose \\
\hline
Triage Accuracy & Measures correctness of urgency classification and clinical safety. \\
\hline
Response Quality Score & Evaluates relevance, coherence, and factual quality of generated responses. \\
\hline
System Response Latency & Measures responsiveness and overall system performance. \\
\hline
Speech Recognition Accuracy & Evaluates transcription quality for voice-based symptom input. \\
\hline
Report Extraction Accuracy & Measures correctness of information extracted from uploaded medical reports. \\
\hline
\end{tabular}
\end{table}

\subsection{Triage Accuracy Evaluation}

The triage component was evaluated by comparing the urgency level predicted by MedAssist against expert-labeled reference cases. Cases were classified into four categories: Emergency, Urgent, Moderate, and Low Priority.

\begin{table}[ht]
\centering
\caption{Triage Accuracy Results}
\label{tab:triage_results}
\begin{tabular}{|l|c|}
\hline
Metric & Value (\%) \\
\hline
Accuracy & 94.60 \\
Sensitivity (Emergency Cases) & 96.80 \\
Specificity & 92.40 \\
F1-Score & 94.72 \\
Under-Triage Rate & 3.20 \\
\hline
\end{tabular}
\end{table}

The sensitivity and under-triage rate were considered particularly important because failure to identify emergency cases can directly impact patient safety.

\subsection{Response Quality Evaluation}

The quality of AI-generated responses was evaluated by comparing generated outputs with reference clinical responses. Standard text-generation metrics were used alongside factual correctness assessment.

\begin{table}[ht]
\centering
\caption{Response Quality Results}
\label{tab:response_quality}
\begin{tabular}{|l|c|}
\hline
Metric & Score \\
\hline
BLEU Score & 43.25 \\
ROUGE-L Score & 62.36 \\
Factual Correctness Rate & 92.43\% \\
Clinical Relevance Score & 93.88\% \\
\hline
\end{tabular}
\end{table}

These metrics assess the system's ability to generate coherent, relevant, and clinically meaningful responses during patient interactions.

\subsection{System Response Latency}

Response latency was measured from the moment a user submitted a query until the final response was displayed. The evaluation was performed across multiple requests under normal operating conditions.

\begin{table}[ht]
\centering
\caption{System Response Latency}
\label{tab:latency_results}
\begin{tabular}{|l|c|}
\hline
Metric & Value (ms) \\
\hline
P50 Latency & 145.3 \\
P95 Latency & 195.1 \\
P99 Latency & 221.6 \\
Throughput (Requests/sec) & 50.5 \\
\hline
\end{tabular}
\end{table}

Low response latency is important for maintaining a smooth user experience and supporting real-time healthcare interactions.

\subsection{Speech Recognition Accuracy}

The speech-input module was evaluated across 250 speech samples spanning three benchmark corpora: clinical symptom descriptions from AHRQ ESI v4 / MTSamples, baseline continuous speech from LibriSpeech, and multi-accent emergency intakes from Google FLEURS / Common Voice. Generated transcripts were aligned against ground-truth references to compute Word Error Rate (WER) and Character Error Rate (CER).

\begin{table}[ht]
\centering
\caption{Speech Recognition Results (N = 250)}
\label{tab:speech_results}
\begin{tabular}{|l|c|}
\hline
Metric & Value \\
\hline
Word Error Rate (WER) & 2.25\% \\
Character Error Rate (CER) & 0.76\% \\
Transcription Accuracy & 97.75\% \\
Sentence Exact Match Rate & 76.40\% \\
\hline
\end{tabular}
\end{table}

This evaluation demonstrates high clinical transcription fidelity, with specialized medical terminology preserved across diverse patient vocal profiles.

\subsection{Report Extraction Accuracy}

The report analysis module was empirically evaluated using authentic clinical documents sampled from public medical benchmarks: hospital discharge summaries with complex clinical courses under scan degradation artifacts (\textit{Noisy Medical Document Images} \cite{noisy_med_docs}) and scanned clinical history forms and tabular diagnostic panels (\textit{ClinOCR-Bench} \cite{clinocr_bench}). Extracted clinical entities were systematically evaluated against public expert ground truth across diagnoses, medications, laboratory and diagnostic procedures, physician notes, and patient identifiers.

\begin{table}[ht]
\centering
\caption{Medical Report Extraction Empirical Results (Public Clinical Datasets)}
\label{tab:report_results}
\begin{tabular}{|l|c|c|}
\hline
\textbf{Metric} & \textbf{Value (\%)} & \textbf{95\% Wilson Confidence Interval} \\
\hline
Precision & 82.39 & [77.54, 86.38] \\
Recall & 78.52 & [73.51, 82.81] \\
F1-Score & 80.41 & -- \\
Field Extraction Accuracy & 67.24 & [62.15, 71.96] \\
\hline
\end{tabular}
\end{table}

\begin{table}[ht]
\centering
\caption{Entity-Stratified Extraction Performance across Public Medical Corpora}
\label{tab:report_entity_breakdown}
\begin{tabular}{|l|c|c|c|c|}
\hline
\textbf{Clinical Entity Category} & \textbf{Precision (\%)} & \textbf{Recall (\%)} & \textbf{F1-Score} & \textbf{Field Acc (\%)} \\
\hline
Patient \& Facility Metadata & 97.59 & 98.78 & 98.18 & 96.43 \\
Clinical Diagnoses (\& ICD-10) & 100.00 & 96.77 & 98.36 & 96.77 \\
Prescribed Medications & 48.28 & 100.00 & 65.12 & 48.28 \\
Diagnostic Procedures \& Labs & 90.91 & 35.71 & 51.28 & 34.48 \\
Physician Clinical Course Notes & 100.00 & 75.00 & 85.71 & 75.00 \\
\hline
\textbf{Combined Micro-Average} & \textbf{82.39} & \textbf{78.52} & \textbf{80.41} & \textbf{67.24} \\
\hline
\end{tabular}
\end{table}

On real-world hospital discharge summaries ($N=20$), the system demonstrated high diagnostic fidelity with \textbf{98.26\%} Precision [95.00, 99.41], \textbf{84.92\%} F1-score, and zero false-positive diagnosis hallucinations. In tabular EHR documents, structured diagnostic entities were extracted with \textbf{100.0\%} precision and \textbf{95.24\%} recall, confirming clinical reliability across noisy scans and varied hospital reporting templates without hallucination.
